% Include Preamble 
\documentclass[10pt,dvipsnames, aspectratio=169]{beamer}
\usetheme[progressbar=frametitle]{metropolis}
\usepackage[]{verbatim}\usepackage[]{}
\usepackage{booktabs}
\usepackage[scale=2]{ccicons}
\usepackage[misc]{ifsym}
\usepackage{wasysym}
\usepackage{listings}
\usepackage{xspace}
\usepackage{amsmath}
\usepackage{xcolor}
\usepackage{ragged2e}\justifying % for justify content % Customize 
\setlength{\parskip}{5pt} % vertical spacing between 2 paragraphs
\setbeamersize{text margin left=12mm, text margin right=12mm} 
\setbeamertemplate{frametitle}[default][left, leftskip=8mm]
%\usepackage[hidelink]{hyperref}
\newcommand{\themename}{\textbf{\textsc{metropolis}}\xspace}

% WARNING: Don't Touch,if you are not compfortable! 
\addtobeamertemplate{frametitle}{}{
%\begin{tikzpicture}[remember picture,overlay]
%\node[anchor=north east,yshift=2pt,xshift=-65pt] at (current page.north east) 
%{\includegraphics[height=.75cm]{img/elixir-portugal-white}};
%\node[anchor=north east,yshift=4pt] at (current page.north east) 
%{\includegraphics[height=1cm]{img/hdro}};
%\end{tikzpicture}
}

% Title Page of Presentation 
% WARNING: Don't change logo, just write your details 
\title{ISCB20.01 - Introduction to LINUX for Biologists}
\date{\today}
\author{Md. Jubayer Hossain\\
        https://jhossain.me/}
\institute{Lead Organizer, Introduction to Scientific Computing for Biologists 
\\ Founder, Health Data Research Organization}
\titlegraphic{\vspace{4cm}\hfill\includegraphics[height=1cm]{img/hdro2}}
\begin{document}
\maketitle
% Sectuion Title 
\section{Section--II: Basic Operations in LINUX	}
% Slide-0
\begin{frame}[t]{Basics}
	\begin{itemize}
		\item  Remember the UNIX/LINUX command line is case sensitive!
		\item The hash (pound) sign \# indicates end of a command and 
		the start of a comment.
		\item The notation \textless \dots \textgreater refers to variables 
		and file names that need to be specified by the user. The symbols 
		\textless  and \textgreater need to be excluded.
	\end{itemize}
\end{frame}


% Slide-0
\begin{frame}[t]{Navigation}\hfill
\begin{center}
	\begin{tabular}{lll} 
		\textbf{Command} & \textbf{Function} & \textbf{Syntax} \\ 
		\hline
		pwd & Present working directory & \color{red} \textdollar pwd \\ 
		ls & List contents of pwd &  \color{red} \textdollar ls  \\ 
		ls -l & Similar as ls, but provides additional info &  \color{red} 
		\textdollar ls -l \\ 
		ls -la & Includes hidden files (.name) as well &  \color{red} 
		\textdollar ls -la \\ 
	
	\end{tabular}
\end{center}
\end{frame}



% Slide-1
\begin{frame}[t]{Directory Operations}\hfill
	\begin{center}
		\begin{tabular}{lll} 
			\textbf{Command} & \textbf{Function} & \textbf{Syntax} \\ 
			\hline
			mkdir & Creates specified directory & \color{red} \textdollar mkdir 
			dir\textunderscore name\\ 
			cp &Copy file/directory as specified in path (-r to include content 
			in directories) & \color{red} \textdollar mkdir 
			dir\textunderscore name\\ 
			
			
			
			
			rmdir & Removes empty directory & \color{red}  \textdollar rmdir 
			dir\textunderscore name	\\  
			rmdir -r & Removes directory including its content & \color{red} 
			\textdollar rm 
			-r 	dir\textunderscore name	\\  
			rmdir -rf & Removes directory including its content & \color{red} 
			\textdollar rm -rf 	dir\textunderscore name
		\end{tabular}
	\end{center}
\end{frame}



% Slide-2
\begin{frame}[t]{File Operations}\hfill
	\begin{center}
		\begin{tabular}{lll} 
			\textbf{Command} & \textbf{Function} & \textbf{Syntax} \\ 
			\hline
			touch & Creates specified file & \color{red} \textdollar touch 
			file\textunderscore name\\ 
			rm & Removes empty file & \color{red}  \textdollar rm 
			file\textunderscore name	\\  
			rm -r & Removes file including its content & \color{red} 
			\textdollar rm -r 	file\textunderscore name	\\  
			rmdir -rf & Removes directory including its content & \color{red} 
			\textdollar rm -rf 	file\textunderscore name
		\end{tabular}
	\end{center}
\end{frame}



% Slide-3
\begin{frame}[t]{File Manipulations}\hfill
	\begin{center}
		\begin{tabular}{lll} 
			\textbf{Command} & \textbf{Function} & \textbf{Syntax} \\ 
			\hline
			touch & Create files & touch \color{red} filenmae.extension \\ 
			pwd & Present working directory & pwd \\ 
			cd & Change directory & \color{red}cd [Directory] \\ 
			cd .. & Previous Directory & \color{red} cd ..
		\end{tabular}
	\end{center}
\end{frame}



% Slide-4
\begin{frame}[t]{Compression}\hfill
	\begin{center}
		\begin{tabular}{lll} 
			\textbf{Command} & \textbf{Function} & \textbf{Syntax} \\ 
			\hline
			touch & Create files & touch \color{red} filenmae.extension \\ 
			pwd & Present working directory & pwd \\ 
			cd & Change directory & \color{red}cd [Directory] \\ 
			cd .. & Previous Directory & \color{red} cd ..
		\end{tabular}
	\end{center}
\end{frame}






% Slide-5
\begin{frame}[t]{Shortcuts}\hfill
	\begin{center}
		\begin{tabular}{lll} 
			\textbf{Command} & \textbf{Function} & \textbf{Syntax} \\ 
			\hline
			touch & Create files & touch \color{red} filenmae.extension \\ 
			pwd & Present working directory & pwd \\ 
			cd & Change directory & \color{red}cd [Directory] \\ 
			cd .. & Previous Directory & \color{red} cd ..
		\end{tabular}
	\end{center}
\end{frame}




% Slide-end
\begin{frame}[t]{References}
	
\end{frame}

% Thank you slide 
\plain{Thank You\\ \ \\ \Huge{\smiley}}

\end{document}